\chapter{Evolución técnica del desarrollo}
\label{anexo:evolucion-tecnica}

Este anexo reúne la reconstrucción cronológica del desarrollo: cómo el proyecto pasó de una colección de cuadernos de experimentación a un producto distribuido en tres repositorios. Se ubica fuera del cuerpo principal porque describe el proceso y no las capacidades del prototipo, que el Capítulo 11 desarrolla. Su lectura resulta útil para comprender el origen de cada componente y la nomenclatura de los bloques de trabajo empleada en el Capítulo 12.

\section{Evolución técnica alcanzada}

Esta sección reconstruye la evolución técnica del producto desde los primeros desarrollos hasta el estado congelado para la entrega del 50 \%. La reconstrucción se apoya en el historial de los tres repositorios del proyecto y en la documentación de checkpoints que cada uno conserva. El criterio expositivo no es cronológico ni cuantitativo: cada etapa se presenta indicando qué limitación resolvía, qué componente la resolvió primero, qué cambio exigió en los repositorios restantes y qué capacidad de producto quedó disponible como resultado.

Antes de recorrer las etapas conviene advertir que la nomenclatura de trabajo no fue uniforme durante todo el desarrollo, y que el documento utiliza tres sistemas de identificación que no deben confundirse entre sí. Esa distinción se detalla en la sección 12.2 y se resume aquí porque condiciona la lectura de toda la sección: las épicas E1 a E10 son unidades académicas derivadas del user research; los identificadores AI-, BE-, DEV-, FE-RD- y FE-P corresponden a backlogs internos de cada repositorio; y los checkpoints P8, P9, P10, P10.5 y siguientes son bloques coordinados de producto.

\subsection{Etapa fundacional y separación en tres repositorios}
\label{anexo:f-etapa-fundacional}

El proyecto no comenzó como un producto distribuido. La primera etapa se concentró en un único repositorio de experimentación, donde el trabajo se organizó como una serie de notebooks numerados que la documentación interna identifica con el prefijo E seguido de un número. Conviene señalar de inmediato que estos hitos E del módulo de IA no guardan relación con las épicas académicas E1 a E10 definidas en el Capítulo 8: son sistemas de numeración independientes que coinciden por casualidad en la letra inicial.

La secuencia experimental avanzó desde la exploración y visualización del dataset y el preprocesamiento y normalización de los volúmenes, hacia un baseline de segmentación sagital que fue sucesivamente diagnosticado, mejorado y validado contra un conjunto de retención. Sobre esa base se abrió un spike axial que exigió inventariar un dataset distinto, resolver el emparejamiento entre imágenes y máscaras, decodificar correctamente las etiquetas y recién entonces entrenar un baseline axial propio. El cierre de esta etapa produjo los dos entrenamientos finales por plano, el mapeo de clases axiales y un primer pipeline de inferencia multiplanar.

La limitación que cerró esta etapa fue de naturaleza arquitectónica, no algorítmica: los modelos funcionaban dentro de notebooks, pero un notebook no es un producto. Los hitos finales de la serie experimental se dedicaron precisamente a resolver ese pasaje, incorporando un orquestador del agente de IA, una traducción del comportamiento hacia un servicio consumible por backend, un smoke de integración, y sendos andamiajes iniciales de servicio Spring Boot y de interfaz de worklist. Esa constatación motivó la separación en tres repositorios independientes (módulo de IA, backend y frontend), con la responsabilidad de cada capa delimitada por contrato.

La separación habilitó el desarrollo en paralelo, pero introdujo el problema que domina el resto de la evolución: mantener coherentes tres repositorios que avanzan a ritmos distintos.

\subsection{Backlogs técnicos por repositorio}
\label{anexo:f-backlogs-repositorios}

Tras la separación, y antes de que existiera coordinación formal entre repositorios, cada uno adoptó su propio sistema de identificadores: AI- en el módulo de inteligencia artificial, BE- en el backend, FE- y FE-RD- en el frontend, y DEV- para la contenedorización y la orquestación comunes a los tres. Ese período explica por qué el producto no apareció terminado alrededor de P9: cada repositorio resolvía su propia progresión antes de que existiera un bloque coordinado.

El módulo de IA avanzó desde la inspección de los checkpoints entrenados hacia el endurecimiento de los manifests con verificación de hash, la carga estricta por plano y la corrida multiplanar con identificador común. El backend partió de congelar los tipos del contrato y avanzó hacia la persistencia verificada sobre PostgreSQL con migraciones versionadas. El frontend construyó los tipos del contrato, el espacio de trabajo por caso, la ejecución, el render de superposiciones y el visor sagital y axial con control de opacidad.

Existe además una serie identificada como FE-P1 a FE-P8, correspondiente a una épica de pulido visual del frontend. Su numeración es interna de ese repositorio y no guarda relación con los checkpoints coordinados de producto: FE-P8 designa un documento de control de calidad de la interfaz, mientras que P8 designa el checkpoint de persistencia.

\subsection{Bloques coordinados de producto}
\label{anexo:f-bloques-coordinados}

% CROSS_REFERENCE_MAPPING_REQUIRED: Tabla 9.1
A partir de P8 el desarrollo se organiza en bloques coordinados entre repositorios. Cada bloque parte de una limitación concreta del producto, se resuelve en el componente que puede resolverla y deja disponible una capacidad nueva junto con la dependencia que abre para el bloque siguiente. La Tabla 9.1 recorre esa secuencia.

\begin{longtable}{|>{\raggedright\arraybackslash}p{0.13\textwidth}|>{\raggedright\arraybackslash}p{0.20\textwidth}|>{\raggedright\arraybackslash}p{0.27\textwidth}|>{\raggedright\arraybackslash}p{0.20\textwidth}|>{\raggedright\arraybackslash}p{0.10\textwidth}|}
\caption{Bloques coordinados de producto, de P8 a P10.9.}
\label{tab:anexo-f-bloques-coordinados}\\
\hline
\textbf{Bloque} & \textbf{Limitación que resolvía} & \textbf{Resolución} & \textbf{Capacidad resultante} & \textbf{Estado} \\
\hline
\endfirsthead
\hline
\textbf{Bloque} & \textbf{Limitación que resolvía} & \textbf{Resolución} & \textbf{Capacidad resultante} & \textbf{Estado} \\
\hline
\endhead
P8 · Persistencia & El sistema funcionaba pero no recordaba: cerrar el prototipo equivalía a perder el trabajo, y la revisión profesional es por definición diferida. & Backend y frontend, sin participación del módulo de IA. Worklist servida desde la base, persistencia de assets y revisiones, historial por referencia de sujeto. El frontend dejó de completar con valores fabricados los campos ausentes. & Reapertura de un estudio procesado con sus resultados, assets y estado de revisión. & Terminado \\
\hline
P9 · Contrato multiplanar & El acoplamiento entre capas se sostenía en acuerdos implícitos: cada cambio en la salida de la inferencia obligaba a corregir en tres lugares. & El módulo de IA normalizó la segunda versión del contrato; el backend lo validó y canonizó en lugar de retransmitirlo; el frontend consumió una forma única. & Espacio de trabajo multiplanar estable. & Terminado \\
\hline
P10 · Seguridad y trazabilidad & Un error interno podía revelar detalles de implementación y no existía forma de reconstruir qué había ocurrido en una corrida fallida. & Autenticación, activación de cuentas, tratamiento de errores, traza distribuida, sanitización de mensajes, auditoría y observabilidad. Contrapartida de sesión en el frontend. & Fronteras explícitas de seguridad y degradación controlada. & Consolidado \\
\hline
P10.5 · Evolución volumétrica & El producto exponía únicamente el corte analizado, cuando el volumen contiene muchos y un resultado sin contexto de vecindad es difícil de evaluar. & Seis bloques encadenados. P10.5-A definió el contrato navegable y es el único cerrado en los tres repositorios a la vez; P10.5-B a P10.5-F distribuyeron catálogo, serving, visor, vinculación por corte y reapertura. & Navegación por cortes con el resultado vinculado a su corte de origen. & Implementado; integración final del frontend pendiente \\
\hline
P10.6 · Estenosis subarticular & — & Modelo entrenado y evaluado internamente sobre RSNA / LumbarDISC, con checkpoint congelado. & Capacidad no disponible de forma automática: el modelo axial no incluye clase de vértebra y falta la generación de la región de interés que el clasificador requiere. & Bloqueado \\
\hline
P10.7 · Hallazgos por nivel & — & Ocho tipos de hallazgo degenerativo, sujetos a revisión profesional y con estado de despliegue asociado. Las probabilidades internas no se exponen. & Salida asistiva por nivel. Los ocho tipos no tienen el mismo soporte empírico y no se presentan como equivalentes. & Implementado \\
\hline
P10.8 · Preflight y guardas & — & Cierre del preflight de expansión clínica e incorporación de guardas, sin entrenamiento adicional. & Verificación de factibilidad previa a la ejecución. & Cerrado \\
\hline
P10.9 · Consolidación & — & Congelamiento del backend y del módulo de IA para el traspaso al frontend, con recorrido técnico verificado mediante un runner automatizado. & Frontera común para continuar el desarrollo. & Congelado; recorrido con interfaz pendiente \\
\hline
\end{longtable}

Sobre P10.5 corresponde una precisión de método: el contrato quedó incorporado a la rama principal de los tres repositorios y la navegación volumétrica está disponible en el frontend, pero la integración definitiva de los contratos congelados en P10.9 permanece pendiente, de modo que el recorrido extremo a extremo con interfaz todavía no se ejecutó.

\subsection{Estado resultante para la entrega del 50 \%}
\label{anexo:f-estado-entrega-50}

El estado del proyecto al cierre de esta versión admite tres categorías que no deben mezclarse.

\textbf{Hecho.} La separación en tres repositorios con responsabilidades delimitadas por contrato; los modelos sagital y axial entrenados, congelados y ejecutables sobre entrada real; el runtime de inferencia con gestión de artefactos verificados por hash; la persistencia de estudios, corridas, assets y revisiones sobre PostgreSQL con migraciones versionadas; la reapertura de estudios ya procesados; la navegación por cortes con el resultado vinculado a su corte de origen; y la revisión profesional obligatoria.

\textbf{Pendiente.} La integración final de los contratos congelados en el frontend; el recorrido extremo a extremo con interfaz; la validación profesional sobre el producto integrado; el despliegue objetivo; y la consolidación de la evidencia final.

\textbf{Bloqueado.} La automatización de P10.6, por ausencia de una región de interés axial automática validada. Es un bloqueo de capacidad y no un trabajo pendiente de ejecución: su resolución depende de una etapa de localización anatómica que no está disponible y cuyos sustitutos se descartaron por razones metodológicas.

La tabla siguiente resume la correlación entre repositorios a lo largo de las etapas descritas, indicando en cada caso qué capacidad de producto quedó disponible y qué dependencia la habilitó.

\section{Evolución del producto visible}
\label{anexo:f-evolucion-producto-visible}

La primera versión del proyecto se concentró en demostrar la viabilidad técnica de un flujo sagital: recibir un volumen de resonancia magnética lumbar, seleccionar un corte representativo, ejecutar un modelo de segmentación 2D y presentar una salida visual revisable. Este recorte permitió validar el pipeline mínimo de datos, inferencia y visualización sin incorporar desde el inicio todas las funciones de una plataforma. Las capacidades que se describen en el resto del capítulo no provienen de un único checkpoint: se acumularon a lo largo de la evolución reconstruida en este anexo, que separa responsabilidades entre repositorios, incorpora el plano axial, consolida la persistencia, normaliza el contrato multiplanar, endurece seguridad y trazabilidad, y evoluciona hacia la navegación volumétrica antes de llegar al tramo P10.6 a P10.9.

El user research modificó ese alcance inicial. Los profesionales entrevistados señalaron que la revisión de una RM lumbar no se limita a un único plano y que los cortes axiales aportan información relevante para la lateralidad, la relación con raíces y la correlación de hallazgos. A partir de esa evidencia se incorporó un segundo módulo de inferencia axial y se reorganizó la interfaz alrededor de una corrida multiplanar.

La evolución del producto puede resumirse en los siguientes hitos:

\begin{itemize}
  \item MVP sagital con inferencia y salida visual básica.
  \item Incorporación del Backend como orquestador y punto único de acceso.
  \item Separación del AI Module como servicio de inferencia independiente.
  \item Persistencia de usuarios, estudios, corridas, resultados y revisiones.
  \item Incorporación del modelo axial y coordinación sagital–axial.
  \item Reapertura de resultados sin repetir la inferencia.
  \item Autenticación, roles y autorización por recurso.
  \item Workspace multiplanar con overlays, mediciones y revisión profesional.
  \item Definición del contrato volumétrico P10.5-A.
  \item Ingesta de estudio DICOM completo y procesamiento full-series (Backend + AI Module), línea de hallazgos degenerativos por nivel (P10.7) y visor navegable implementado; Backend y AI Module congelados para el traspaso al Frontend (P10.9), con la integración final del contrato en el Frontend pendiente.
\end{itemize}

La incorporación axial no reemplaza el modelo sagital ni implica que ambos modelos segmenten las mismas clases. Cada plano conserva su dataset, preprocesamiento, artifact, clases de salida y criterios de validación. El prototipo integra los resultados en una experiencia común, pero no fusiona artificialmente modelos incompatibles.

\section{Evolución del pipeline de inferencia}
\label{anexo:f-evolucion-pipeline-inferencia}

El desarrollo se organizó en dos ciclos complementarios que, vistos desde este componente, recorren la cadena experimentación en notebooks, backlog técnico propio con prefijo AI-, runtime real, normalización del contrato en P9, endurecimiento en P10, evolución volumétrica en P10.5 y extensión en el tramo P10.6 a P10.9; la correlación con los demás repositorios se desarrolla en este anexo y no se repite aquí. El primero resolvió incertidumbres sobre datasets, máscaras, taxonomías, orientación, preprocesamiento y entrenamiento. El segundo convirtió ese conocimiento en un producto integrado y, en el tramo P10.6–P10.9, incorporó una línea de hallazgos degenerativos asistivos por nivel (P10.7) a partir de Sagital T1 y Sagital T2 independientes, un modelo de estenosis subarticular en Axial T2 evaluado internamente pero aún sin ROI automática (P10.6), el cierre de preflight y guardas sin entrenamiento adicional (P10.8) y la consolidación y congelamiento del componente para el traspaso al Frontend (P10.9).

\begin{longtable}{|>{\raggedright\arraybackslash}p{0.20\textwidth}|>{\raggedright\arraybackslash}p{0.31\textwidth}|>{\raggedright\arraybackslash}p{0.25\textwidth}|>{\raggedright\arraybackslash}p{0.16\textwidth}|}
\caption{Evolución de la implementación del pipeline de IA.}
\label{tab:anexo-f-pipeline-inferencia}\\
\hline
\textbf{Etapa} & \textbf{Resultado de implementación} & \textbf{Evidencia principal} & \textbf{Estado al 50 \%} \\
\hline
\endfirsthead
\hline
\textbf{Etapa} & \textbf{Resultado de implementación} & \textbf{Evidencia principal} & \textbf{Estado al 50 \%} \\
\hline
\endhead
Exploración y preparación & Comprensión de formatos, máscaras, clases y orientación de los datasets. & Notebooks, inventarios, figuras de control y decisiones de pairing. & Consolidado experimentalmente. \\
\hline
Entrenamiento y consolidación & Modelos 2D multiclase diferenciados por plano. & Checkpoints E5, E10 y E12; métricas por clase. & Consolidado experimentalmente. \\
\hline
Registro final & Nombres estables, modelKey, manifests y artifacts finales identificados. & Directorio models/final y notebook de registro final. & Consolidado. \\
\hline
Productización & Lógica experimental trasladada a módulos reutilizables del servicio. & Registry, materializador, arquitecturas, runtime y API FastAPI. & Implementado. \\
\hline
Ejecución real & Carga de state\_dict, model.eval(), torch.inference\_mode() y generación de outputs. & Runtime PyTorch y pruebas automatizadas. & Implementado; evidencia final se consolida en el Capítulo 13. \\
\hline
Orquestación multiplanar & Corridas sagital y axial coordinadas bajo un contexto común. & run\_multiplanar\_pipeline y endpoint multiplanar. & Implementado. \\
\hline
Evolución volumétrica & Catálogo de series, serving de cortes y procesamiento full-series integrados en Backend, AI Module y visor. & P10.5-A y fixtures compartidos. & Implementado (full-series). \\
\hline
\end{longtable}

\section{Organización interna del módulo de inteligencia artificial}
\label{anexo:f-organizacion-interna-ai}

El módulo de inteligencia artificial agrupa responsabilidades que no corresponden al backend de producto, de modo que el código de inferencia pueda probarse sin depender de la interfaz y que la identidad del modelo no quede implícita en rutas o convenciones manuales. Un registro de modelos resuelve la clave, la versión, las clases y la dimensión de entrada de cada artifact; un materializador lo descarga desde el almacenamiento externo mediante archivo temporal y reemplazo atómico, de modo que nunca queda un artifact parcial; el runtime resuelve el dispositivo, carga y mantiene en caché los modelos y ejecuta la inferencia; un orquestador agrupa ambos planos bajo identificadores comunes sin fusionar sus taxonomías; y una capa de servicio expone los accesos que el backend consume.

La reconstrucción de las arquitecturas se realiza con carga estricta: si los nombres de capa del artifact no coinciden con los de la arquitectura declarada, la carga falla de forma explícita en lugar de completar los pesos faltantes con valores iniciales. Cada artifact conserva un manifest con sus clases, su configuración, sus métricas y su resumen criptográfico, y el sistema verifica esa correspondencia antes de habilitar el modelo.

La condición de disponibilidad se comprueba antes de aceptar una solicitud: deben existir el artifact local completo, un manifest coherente y una arquitectura compatible. Si alguna de las tres falta, el servicio lo informa como estado de indisponibilidad y no como resultado de análisis.
